\chapter{GPT Model completion and learning loop}
\label{ch:gpt_training}

In this chapter, we build a complete\keyterm{GPT model}by stacking$N$layers of the transformer blocks from Chapter 6, and build a loop to learn this model. The dataset, loss function, optimizer, and scheduler are all implemented directly.

\section{GPT model structure}

GPT is an autoregressive language model with$N$layers of transformer decoder. The overall structure is as follows:

\begin{enumerate}
\item Token ID$\to$[Token + Position Embedding]
  \item [Transformer Block] $\times N$
  \item [Final LayerNorm]
  \item [LM Head: $d\_{model} \to V$]
  \item logits
\end{enumerate}

\begin{figure}[htbp]
\centering
\begin{tikzpicture}[
  every node/.style={font=\small\sffamily},
  box/.style={draw=inkblue, fill=inkblue!8, rounded corners=2pt, minimum width=4cm, minimum height=0.7cm},
  rep/.style={draw=inkgray, fill=inkgray!10, rounded corners=2pt, minimum width=4cm, minimum height=0.7cm, dashed}
]
  \node[box] (in) at (0, 0) {token IDs [batch, seq]};
  \node[box, below=0.4cm of in] (emb) {Embedding (token + position)};
  \node[rep, below=0.4cm of emb] (blk) {Transformer Block $\times N$};
  \node[box, below=0.4cm of blk] (ln) {Final LayerNorm};
  \node[box, below=0.4cm of ln] (head) {LM Head ($d\_{model} \to V$)};
  \node[box, below=0.4cm of head] (out) {logits [batch, seq, V]};

  \draw[-{Stealth[length=4pt]}, inkblue] (in) -- (emb);
  \draw[-{Stealth[length=4pt]}, inkblue] (emb) -- (blk);
  \draw[-{Stealth[length=4pt]}, inkblue] (blk) -- (ln);
  \draw[-{Stealth[length=4pt]}, inkblue] (ln) -- (head);
  \draw[-{Stealth[length=4pt]}, inkblue] (head) -- (out);

  % Weight tying indication
  \draw[inkred, dashed, <->] (emb.east) to[bend left=30] node[right, inkred, font=\scriptsize] {weight tying} (head.east);
\end{tikzpicture}
\caption{GPT model structure. Embedding weights and LM Head Share weights (weight tying) to save parameters.}
\label{fig:gpt}
\end{figure}

\section{weighted tie (Weight Tying)}

The embedding matrix$E$($V \times d\_{model}$) and the LM Head matrix$W\_{out}$($d\_{model} \times V$) actually contain the same information. The embedding of token$i$is$E[i]$, and$W\_{out}^T[i]$is used when calculating logit. Sharing these two can save the number of parameters by$V \times d\_{model}$.

\begin{lstlisting}[language=Python]
class GPT(nn.Module):
    def __init__(self, config: ModelConfig):
        super().__init__()
        self.config = config
        self.embedding = EmbeddingStack(
            vocab_size=config.vocab_size,
            d_model=config.d_model,
            max_len=config.context_length,
            pad_id=config.pad_token_id,
        )
        self.blocks = nn.ModuleList([
            TransformerBlock(config, layer_norm_style="pre")
            for _ in range(config.n_layers)
        ])
        self.ln_f = nn.LayerNorm(config.d_model, eps=config.layer_norm_eps)
        self.lm_head = nn.Linear(config.d_model, config.vocab_size, bias=False)
        # weighted tie
        self.lm_head.weight = self.embedding.token_emb.embedding.weight
\end{lstlisting}

\begin{notebox}[Advantages and disadvantages of weighted ties]
\textbf{merit}: Parameter saving (about 1/3 reduction based on GPT-2 small), consistency of embedding and output.

\textbf{disadvantage}: Embedding also receives the gradient of the output layer, making learning more complicated. Some models (such as LLaMA) are trained separately without using ties.

GPT-2 and GPT-3 use ties. LLaMA and Mistral are not used. This book uses ties following the GPT style.
\end{notebox}

\section{Forward Pass}

\begin{lstlisting}[language=Python]
    def forward(self, token_ids, return_hidden=False):
        batch, seq = token_ids.shape
        # Create a causal mask
        mask = make_causal_mask(seq, device=token_ids.device)
        # Embedding
        x = self.embedding(token_ids)
        # Transformer block pass
        for block in self.blocks:
            x = block(x, mask=mask)
        # Final normalization (Pre-LNso it is needed at the end)
        x = self.ln_f(x)
        # logit
        logits = self.lm_head(x)
        if return_hidden:
            return logits, x
        return logits
\end{lstlisting}

\begin{infobox}[Why do you need LayerNorm at the end?]
The Pre-LN block normalizes the input of each block. However, the output of the last block is added to the residual stream in an unnormalized state. This must be normalized before inserting it into the LM Head to ensure a stable output distribution. So\code{ln\_f}(final LayerNorm) is needed.

In Post-LN, each block output is already normalized, so\code{ln\_f}is unnecessary.
\end{infobox}

\section{dataset: sliding windows}

Language modeling datasets tokenize text and then split it into fixed-length chunks. In each chunk, input is the first$L-1$token, target is the last$L-1$token (shift one space).

\begin{lstlisting}[language=Python]
class LMDataset(Dataset):
    def __init__(self, text, tokenizer, context_length=128,
                 add_bos=True, add_eos=True):
        self.context_length = context_length
        # Tokenize text into long sequences
        if isinstance(text, str):
            text = [text]
        all_ids = []
        for t in text:
            ids = tokenizer.encode(t, add_bos=add_bos, add_eos=add_eos)
            all_ids.extend(ids)
        self._ids = all_ids
        # context_length+Split into 1 size chunks
        cl = context_length + 1
        self._chunks = []
        for i in range(0, len(all_ids), cl):
            chunk = all_ids[i:i+cl]
            if len(chunk) < cl:
                chunk = chunk + [tokenizer.special.pad_id] * (cl - len(chunk))
            self._chunks.append(chunk)

    def __getitem__(self, idx):
        chunk = self._chunks[idx]
        input_ids = torch.tensor(chunk[:-1], dtype=torch.long)
        target_ids = torch.tensor(chunk[1:], dtype=torch.long)
        return input_ids, target_ids
\end{lstlisting}

\begin{figure}[htbp]
\centering
\begin{tikzpicture}[
  every node/.style={font=\small\sffamily},
  tok/.style={draw=inkblue, fill=inkblue!15, minimum width=0.5cm, minimum height=0.5cm},
  dstin/.style={draw=inkblue, fill=inkblue!30, minimum width=0.5cm, minimum height=0.5cm},
  tgt/.style={draw=inkred, fill=inkred!30, minimum width=0.5cm, minimum height=0.5cm}
]
  % sequence
  \node[inkgray, anchor=east] at (-0.3, 0) {tokens:};
  \foreach \i/\t in {0/B, 1/t, 2/h, 3/e, 4/c, 5/a, 6/t, 7/E} {
    \node[tok] at (\i*0.55, 0) {\t};
  }

  \node[inkgray, anchor=east] at (-0.3, -0.8) {input:};
  \foreach \i/\t in {0/B, 1/t, 2/h, 3/e, 4/c, 5/a, 6/t} {
    \node[dstin] at (\i*0.55, -0.8) {\t};
  }

  \node[inkgray, anchor=east] at (-0.3, -1.6) {target:};
  \foreach \i/\t in {0/t, 1/h, 2/e, 3/c, 4/a, 5/t, 6/E} {
    \node[tgt] at (\i*0.55, -1.6) {\t};
  }

  % arrow
  \foreach \i in {0, 1, 2, 3, 4, 5, 6} {
    \draw[-{Stealth[length=3pt]}, inkgray!50] (\i*0.55, -0.3) -- (\i*0.55, -0.55);
  }

  \node[inkgray, font=\scriptsize, anchor=west] at (5, -0.8) {input is one space before target};
  \node[inkgray, font=\scriptsize, anchor=west] at (5, -1.6) {B=<bos>, E=<eos>};
\end{tikzpicture}
\caption{Sliding Window Dataset. inputclass targetis the same sequence one space shiftwhich. The model is inputlooking at targetpredict.}
\label{fig:sliding-window}
\end{figure}

\section{loss function: Cross-Entropy}

The loss of the language model is\keyterm{Cross entropy (cross-entropy)}. Measures the difference between the probability distribution predicted by the model and the correct answer (one-hot).

\begin{formula}[Cross-Entropy Loss]
\[
\mathcal{L} = -\frac{1}{N} \sum\_{i=1}^{N} \log P(y\_i \mid x\_{<i})
\]
$y\_i$is the correct answer token,$P$is the model’s prediction probability. The lower the better.
\end{formula}

\begin{lstlisting}[language=Python]
def cross_entropy_loss(logits, targets, ignore_index=-100):
    # logits: [batch, seq, V], targets: [batch, seq]
    batch, seq, vocab = logits.shape
    logits_flat = logits.view(-1, vocab)
    targets_flat = targets.view(-1)
    return F.cross_entropy(logits_flat, targets_flat, ignore_index=ignore_index)
\end{lstlisting}

\subsection{Label Smoothing (select)}

A technique that improves generalization performance by allowing some uncertainty instead of being 100% sure of the correct answer.

\begin{lstlisting}[language=Python]
def label_smoothing_loss(logits, targets, smoothing=0.1, ignore_index=-100):
    """1 for the correct answer token-smoothing probability, to the rest smoothing/vocab probability."""
    log_probs = F.log_softmax(logits.view(-1, logits.size(-1)), dim=-1)
    targets_flat = targets.view(-1)
    nll_loss = -log_probs.gather(1, targets_flat.clamp(min=0).unsqueeze(1)).squeeze(1)
    smooth_loss = -log_probs.mean(dim=-1)
    mask = (targets_flat != ignore_index).float()
    nll_loss = (nll_loss * mask).sum() / mask.sum().clamp(min=1)
    smooth_loss = (smooth_loss * mask).sum() / mask.sum().clamp(min=1)
    return (1.0 - smoothing) * nll_loss + smoothing * smooth_loss
\end{lstlisting}

\section{optimizer: AdamW}

\keyterm{AdamW}is an optimization algorithm that combines Adam with weight decay. The key is not to apply weight decay to bias and LayerNorm.

\begin{lstlisting}[language=Python]
def build_optimizer(model, config):
    # Separate parameter groups: decay / no_decay
    decay, no_decay = [], []
    for name, p in model.named_parameters():
        if not p.requires_grad: continue
        if p.ndim < 2 or "bias" in name or "ln" in name.lower():
            no_decay.append(p)
        else:
            decay.append(p)
    return AdamW(
        [{"params": decay, "weight_decay": config.weight_decay},
         {"params": no_decay, "weight_decay": 0.0}],
        lr=config.learning_rate, betas=(0.9, 0.95)
    )
\end{lstlisting}

\begin{notebox}[Why are bias and LayerNorm subtracted from weight decay?]
Weight decay is a regularization that pulls the weights towards 0. Matrix weights have the effect of preventing overfitting, but the$\gamma, \beta$of bias and LayerNorm have low dimensions, so the effect of weight decay is minimal and may actually hinder learning. This is standard practice in the PyTorch community.
\end{notebox}

\section{learning rate scheduler}

At the beginning of learning, the learning rate is gradually warmed up and then decreased to a cosine curve. This\keyterm{warmup + cosine}schedule is used as standard in GPT-3.

\begin{lstlisting}[language=Python]
def build_scheduler(optimizer, config, total_steps):
    warmup = config.warmup_steps
    def lr_lambda(step):
        if step < warmup:
            return float(step) / max(1, warmup)
        progress = (step - warmup) / max(1, total_steps - warmup)
        if config.lr_scheduler == "cosine":
            return 0.5 * (1.0 + math.cos(math.pi * min(progress, 1.0)))
        return 1.0
    return LambdaLR(optimizer, lr_lambda)
\end{lstlisting}

\begin{figure}[htbp]
\centering
\begin{tikzpicture}[
  scale=0.9,
  every node/.style={font=\small\sffamily}
]
  % axis
  \draw[inkblue, line width=0.5pt, ->] (0, 0) -- (11, 0) node[right] {step};
  \draw[inkblue, line width=0.5pt, ->] (0, 0) -- (0, 3.5) node[above] {LR};
  % warmup section (linear increase)
  \draw[inkred, line width=1pt] (0, 0) -- (2, 2.5);
  \node[inkred, font=\scriptsize, anchor=south] at (1, 2.6) {warmup};
  % cosine section
  \draw[inkblue, line width=1pt, domain=2:10, samples=50] plot (\x, {2.5*0.5*(1+cos(180*(\x-2)/8))});
  \node[inkblue, font=\scriptsize, anchor=south] at (6, 2.6) {cosine decay};
  % Maximum LR displayed
  \draw[inkgray, dashed] (0, 2.5) -- (2, 2.5);
  \node[inkgray, font=\scriptsize, anchor=east] at (-0.1, 2.5) {$lr\_{max}$};
\end{tikzpicture}
\caption{Warmup + Cosine Learning rate schedule. At first, raise it linearly., It then decreases to a cosine curve..}
\label{fig:lr-schedule}
\end{figure}

\subsection{why warmupIs this necessary??}

At the beginning of training, the weights are close to random. If a large learning rate is applied immediately, the gradient may become unstable and diverge. If you start with a small learning rate and increase it gradually, the model enters a stable region.

In particular, the initial gradient of attention can be very large (softmax saturation, etc.), and warmup alleviates this. Pre-LN structures are less sensitive to warm-up than post-LN, but are still safe to use.

\section{gradient clipping}

Limit the gradient norm to prevent the gradient from being too large and causing the weight to change significantly.

\begin{lstlisting}[language=Python]
if config.grad_clip > 0:
    torch.nn.utils.clip_grad_norm_(model.parameters(), config.grad_clip)
\end{lstlisting}

Usually\code{grad\_clip = 1.0}is used. If the gradient norm exceeds 1.0, it is scaled to 1.0.

\section{trainer}

Create a learning loop by tying all components together.

\begin{lstlisting}[language=Python]
class Trainer:
    def __init__(self, model, train_dataset, config, model_config, eval_dataset=None):
        self.model = model
        self.config = config
        self.device = torch.device(config.device)
        self.model.to(self.device)
        self.train_loader = DataLoader(train_dataset, batch_size=config.batch_size,
                                       shuffle=True, collate_fn=collate_batch)
        # Calculate total number of steps
        total_steps = len(train_dataset) // config.batch_size * config.max_epochs
        self.optimizer = build_optimizer(model, config)
        self.scheduler = build_scheduler(self.optimizer, config, total_steps)
        # ...

    def train(self):
        self.model.train()
        for inputs, targets in self.train_loader:
            inputs, targets = inputs.to(self.device), targets.to(self.device)
            logits = self.model(inputs)
            loss = cross_entropy_loss(logits, targets,
                                       ignore_index=self.model_config.pad_token_id)
            self.optimizer.zero_grad(set_to_none=True)
            loss.backward()
            if self.config.grad_clip > 0:
                torch.nn.utils.clip_grad_norm_(self.model.parameters(),
                                                self.config.grad_clip)
            self.optimizer.step()
            self.scheduler.step()
\end{lstlisting}

\section{Perplexity}

\keyterm{Perplexity (Perplexity, PPL)}is a representative evaluation index of language models.

\begin{formula}[Perplexity]
\[
\text{PPL} = \exp\left(\frac{1}{N}\sum\_i \log P(y\_i | x\_{<i})\right) = \exp(\text{loss})
\]
PPL indicates ``how many candidates, on average, does the model wander through'' when predicting the next token. PPL=1 means perfect, PPL=10 means you can narrow it down to 10 candidates.
\end{formula}

\begin{infobox}[PPL interpretation]
\begin{itemize}
\item PPL = 1: Model is 100% sure of the next token. Theoretical minimum.
\item PPL = 10: Narrow down to one of 10 tokens. A level that is useful as a language model.
\item PPL = 100: One of 100 tokens. Almost random.
\item PPL =$V$(vocabulary size): Completely random. Not learned.
\end{itemize}
WikiText-103 PPL of GPT-2 small is about 17. GPT-3 175B is about 11.
\end{infobox}

\section{Save checkpoint/load}

Model weights can be periodically saved during training and resumed after interruption or used for inference.

\begin{lstlisting}[language=Python]
def save_checkpoint(self, filename):
    path = self.out_dir / filename
    torch.save({
        "model_state": self.model.state_dict(),
        "optimizer_state": self.optimizer.state_dict(),
        "scheduler_state": self.scheduler.state_dict(),
        "global_step": self.global_step,
        "model_config": self.model_config.to_dict(),
        "trainer_config": self.config.to_dict(),
    }, path)

def load_checkpoint(self, path):
    ckpt = torch.load(path, map_location=self.device)
    self.model.load_state_dict(ckpt["model_state"])
    self.optimizer.load_state_dict(ckpt["optimizer_state"])
    self.scheduler.load_state_dict(ckpt["scheduler_state"])
    self.global_step = ckpt["global_step"]
\end{lstlisting}

\section{summation}

\begin{itemize}
  \item GPT = Embedding + Transformer Block $\times N$ + LayerNorm + LM Head.
\item Save parameters by sharing LM Head weights with embedding with weight ties.
\item In Pre-LN structure,\code{ln\_f}is required at the end.
\item Language modeling creates (input, target) pairs with a sliding window.
\item Loss is cross-entropy, optimizer is AdamW (bias/LayerNorm excludes weight decay).
\item Warmup + Cosine learning rate schedule + gradient clipping is standard.
\item Evaluate model quality with Perplexity =$\exp(\text{loss})$.
\end{itemize}

\section{practice problems}

\begin{enumerate}
\item How much will the parameters increase in GPT-2 small (V=50257, d=768) if weight ties are not used?
\item How does learning change if the number of warmup steps is set to 0?
\item What is the effect of applying strong label smoothing to 0.3?
\item What can happen when learning without gradient clipping?
\item Can a model with a PPL of 1.0 actually exist? Why?
\end{enumerate}

The next chapter covers how to generate text with the learned model. Greedy, Top-K, and Top-P sampling are all implemented.
